\chapter{绪论}

\section{研究背景与概况}
就目前的高校科研环境来说，项目规模和经费体量已经大到了一个前所未有的地步。然而，这种快速增长也直接暴露了原有管理手段的短板：许多单位还在依赖人工录入或者功能单一的陈旧财务系统。这种模式在处理现在这种多来源、高维度的大量经费数据时，往往显得非常吃力，数据冗余、交互差、分析维度死板等问题几乎是通病。

实际上，数字化转型已经从单纯的口号变成了高校管理的客观刚需。针对这些痛点，尝试利用大数据处理和多维分析技术来设计一套辅助决策系统，就成了提高管理效能的切入点。如果能从项目跑了多少、钱花到了哪类科目、以往的支出规律是什么样的这些维度进行立体化剖析，管理者就能获得更清晰的全局视野。这不仅能看清当下的预算流向，更能基于积累的数据，为之后的额度核定和资源优化提供一个靠谱的参考。开发一套具备这种分析能力的系统，确实是为了解决科研经费管理中那些真实存在的痛点。

\section{国内外研究现状}
从目前高校的财务管理实践来看，预算编制、报销流程以及后期的绩效评价，每一个环节都存在着或多或少的改进空间。许多传统的管理方式仍停留在跨部门的手工协同阶段，这在客观上导致了大家常说的“信息孤岛”。为了改变监管滞后的现状，学术界和管理界都在探讨如何利用信息化手段来实现全过程的留痕与数据驱动决策。

在国内，信息化研究大多聚焦于流程的优化。吴丽杰曾建议在智慧财务的大背景下，把科研、财务和资产这些系统彻底打通，从而实现经费的全程在线管理\cite{Wu2024}。周勇也从制度执行的角度分析过监管的漏洞，认为规范业务流程和增强部门联动才是提高效率的关键\cite{Zhou2010}。虽然这些研究为数字化转型打下了基础，但就目前的成果看，针对实时可视化决策支持的具体实现，讨论得还是相对较少。

说到可视化，ECharts 这种高性能的绘图引擎正变得越来越普及。谢韵佳在做学籍管理研究时就发现，利用可视化模型对辅助决策有显著优势\cite{Xie2019}。王振辉也通过数据验证了直观的图表对提升管理效率的积极作用\cite{Wang2021}。从国际视角看，Gaftandzhieva 认为基于商业智能的监控系统已经成了高校管理的新范式，能显着强化预算执行的力度\cite{Gaftandzhieva2023}。此外，张翠莲设计的交互式查询方案也极大地提升了数据的透明度与易读性\cite{Zhang2025}。

技术架构方面，Spring Boot 凭借稳健的表现得到了广泛认可。相关文献提到，它在逻辑分层和数据持久化上表现出色，适合处理像学生系统这种需要快速迭代的业务\cite{Yang2021, Zhao2020}。但在专门针对科研经费全流程分析的系统化框架方面，目前的研究还是显得比较薄弱。戴建青等人的研究也提到了这点，认为管理视角的碎片化限制了决策的科学性，强调了必须提升数据的集中管理能力\cite{Dai2019}。

此外，预算执行与绩效的关系也备受关注。张川等人的分析表明，目前的监控体系在绩效导向的分析能力上尚不完善\cite{ZhangC2019}。Schiller 的调查揭示了缺乏自动化工具给科研人员带来的沉重负担，认为构建数字化监管体系势在必行\cite{Schiller2022}。基于内控制度的考量，周芩羽强调了通过信息化提升流程透明度的重要性\cite{ZhouQY2020}。综合来看，利用大数据和多维度分析来构建宏观决策体系，已经是当前管理研究的一个主要发展趋势\cite{Xu2023}。

\section{课题要解决的核心点}
在实际操作中，科研经费管理面临的几个问题确实很棘手。这一课题的设计初衷，主要就是想针对下面这三个难点找出解决方案。

关于监控滞后的问题，目前的系统大多偏向于“事后统计”，也就是钱花完了才知道。管理者在项目执行过程中很难及时发现异常。基于这一发现，**我**设计了一个“智能风险标签引擎”，通过对时间进度和报销情况进行实时的多维对比，力求把被动管理转变为主动预警。

决策手段匮乏也是一个急需解决的问题。过去在评估项目或者审批预算调剂时，管理人员手里缺乏直观、多维的数据支撑。本系统通过引入多维分析模型和交互式仪表盘，把复杂的数字直接转化成易于理解的趋势图和对比图，让管理决策变得更加透明。

最后是管理链条没能实现完全闭环。针对预算编制与实际支出脱节的现象，系统强化了从申报、预算审核到报销审计的整体衔接。通过对多源数据的采集，确保每一笔钱的流向都有据可查，从而在很大程度上解决了追踪难的问题。

\section{全文的逻辑框架}
这篇论文主要是围绕这套“决策支持系统”的设计与实现来写的。全书一共分成了七章。

第一章是绪论。主要介绍了选课背景、管理现状以及整篇论文的思路。第二章是技术综述。主要梳理了开发过程中用到的关键技术，包括 Spring Boot、Vue3、MyBatis-Plus 以及绘图用的 ECharts。

第三章是系统需求分析。该部分对业务流程和功能需求做了深度调研，并明确了系统必须具备的多维决策支持能力指标。第四章是总体设计。详细展示了系统的整体架构方案，并在多维建模思路上做了重点探讨。

第五章是核心实现。这章展示了后端服务、预警算法以及多维可视化看板的具体实现细节。第六章是测试与分析。通过功能测试和压力测试来验证系统运行的平稳度，并利用实际数据测试预警的准确度。最后的第七章是总结。归纳了研究工作和创新点，也列出了目前系统存在的不足及以后改进的方向。